% Generated draft: requires mathematical and source review.
\documentclass[11pt]{article}
\usepackage[margin=25mm]{geometry}
\usepackage{fontspec}
\setmainfont{Latin Modern Roman}
\usepackage{amsmath,amssymb,mathtools,unicode-math}
\setmathfont{Latin Modern Math}
\usepackage{enumitem,needspace,xurl,hyperref}
\setlength{\parindent}{0pt}
\setlength{\parskip}{5pt}
\setlength{\emergencystretch}{3em}
\newcommand{\N}{\mathbb{N}}
\newcommand{\Z}{\mathbb{Z}}
\newcommand{\Q}{\mathbb{Q}}
\newcommand{\R}{\mathbb{R}}
\newcommand{\C}{\mathbb{C}}
\newcommand{\F}{\mathbb{F}}
\newcommand{\A}{\mathbb{A}}
\newcommand{\PP}{\mathbb{P}}
\newcommand{\E}{\mathbb{E}}
\newcommand{\eps}{\varepsilon}
\newcommand{\dd}{\,\mathrm{d}}
\begin{document}
\section*{UnsolvedMath ID: 30000208}
\textbf{Source problem code:} catalogue entry
\subsection{Definitions and mathematical statement}
For square matrices of the same size, the matrix principle is
\[
\Phi:\qquad\forall n\geq1\ \forall A,B\in\operatorname{Mat}_n(R),\quad
AB=I_n\ \Longrightarrow\ BA=I_n.
\]
The coefficient structure \(R\), its encoding, and the coding of matrix multiplication must be fixed as in the source; they are not given in the record. The usual principle over a field or commutative ring must not be silently asserted for arbitrary noncommutative rings. Rectangular matrices are not included.

Let \(T_{\mathrm{AC}^0},T_{\mathrm{AC}^0(2)},T_{\mathrm{TC}^0},T_{\mathrm{NC}^1},T_{\mathrm L},T_{\mathrm P}\) denote the particular bounded-arithmetic theories corresponding in the report to the displayed complexity-class hierarchy. The source question is: which is the weakest of these theories that proves the correctly encoded sentence \(\Phi\)? Once the intended theory inclusion ordering is fixed, the target is to identify its least level proving \(\Phi\), rather than to determine the truth of ordinary finite-dimensional linear algebra in an unrestricted background theory. A complexity class alone does not define a unique formal theory.

The original passage reports an upper bound ``at most \(\mathrm P\)'' and conjectures that a lower level suffices. It also mentions associated quantified propositional proof systems and a method translating \(\Sigma^B_1\) theorems into polynomial-size proof families in the associated systems. The precise systems, translation, and formula hypotheses are absent, so this context does not supply an explicit smaller-theory proof or any lower-bound theorem.

The supplied statement verification says the hierarchy was restored from the official report and the OCR-corrupted identity matrix \(U\) was corrected to \(I\). Its literature assessment checked 2026-08-21 still calls the extraction a fragment without a complete theorem to test. The repaired hierarchy gives a definite research direction, but the missing theories, coefficient structure, and encodings still prevent identification of a unique proof-theoretic problem.
\subsection{Short English statement}
Find the least level of the specified bounded-arithmetic hierarchy that proves that a one-sided inverse of a square matrix is also a two-sided inverse. The question concerns formal proof strength; its matrix coefficients and actual formal theories are not specified.
\subsection{Sources}
\begin{itemize}
\item Local UnsolvedMath registry, permanent ID 30000208, catalogue entry, \emph{Proof-Theoretic Strength of a Matrix Principle}.
\item \emph{Mathematical Logic: Proof Theory, Type Theory and Constructive Mathematics} (2005), Oberwolfach Reports, pp. 779--813, DOI \url{https://doi.org/10.4171/owr/2005/14}.
\item Supplied official report link: \url{https://ems.press/content/serial-article-files/45988?nt=1}.
\item Supplied original fragment records the \(\mathrm P\) upper bound and proof-system translation context; supplied source verification restores the hierarchy. Literature triage checked 2026-08-21.
\end{itemize}
\end{document}
